% The cover letter to Cell's editors (build.py letter). Its claims are the paper's own (Summary, Highlights, Introduction,
% Discussion); its numbers come from the ledger. Marked fields are the author's to fill before submission.
\documentclass[11pt, letterpaper]{article}
\usepackage{fontspec}
\IfFontExistsTF{Arial}{\setmainfont{Arial}}{%
  \setmainfont{arial.ttf}[Path=C:/Windows/Fonts/, BoldFont=arialbd.ttf, ItalicFont=ariali.ttf, BoldItalicFont=arialbi.ttf]}
\usepackage[american]{babel}
\usepackage[letterpaper, margin=1in]{geometry}
\usepackage{xcolor}
\usepackage{enumitem}
\usepackage[hyphens]{url}
\usepackage[hidelinks]{hyperref}
\setlength{\parindent}{0pt}
\setlength{\parskip}{0.7\baselineskip}
\setlist{nosep, leftmargin=1.5em, topsep=0pt}
\pagestyle{empty}
% the paper's macros that the letter uses (cellpress.sty)
\newcommand{\cellval}[2]{\expandafter\def\csname cell@val@#1\endcsname{#2}}
\makeatletter
\newcommand{\val}[1]{\ifcsname cell@val@#1\endcsname\csname cell@val@#1\endcsname
  \else\PackageError{cover letter}{No ledger value #1}{Rebuild the ledger and values.tex}\fi}
\makeatother
\newcommand{\tofill}[1]{\textcolor{red!70!black}{[#1]}}
\newcommand{\gene}[1]{\textit{#1}}
\input{generated/values}

\begin{document}
October 25, 2026

Dear Editors,

We are pleased to submit the Article ``Brain state locks silent memories in a trained fly connectome'' for consideration at
\textit{Cell}.

Complete connectomes now map every synapse of an insect's brain and nerve cord, but whether a fixed wiring diagram can learn
demanding behavior, and what its memories consist of, has not been tested: models built on connectomes have been run without
training or trained within one sensory pathway. We trained a spiking model of the entire male \textit{Drosophila} central
nervous system, \val{L01} neurons wired as in the male connectome, to plan moves and to remember a visual cue, and dissected the
memory cell by cell. The study makes three findings.
\begin{itemize}
  \item \textbf{A fixed connectome learns.} Training tunes neuronal parameters and a readout, never the wiring, and the model
  learns both tasks.
  \item \textbf{Working memory is held silently in neuronal adaptation.} Spikes write the cue into adaptation, which holds it
  without extra firing until the next input reads it out. The cells that hold the trace are enriched for transcripts of the
  calcium-activated potassium channel \gene{SK} in fly cell atlases.
  \item \textbf{Recall is locked to the brain state of learning.} The learned readout uses the trace only at the delays it
  learned, because the same trace drives a different response at another delay. A disease-like loss of Kv4 potassium current
  after learning, modeled on the loss that amyloid-β42 causes in flies, rewrites the trace and changes how it is read:
  recall fails, while learning in the new state is spared.
\end{itemize}

These results join three questions that are usually asked apart: what holds an activity-silent memory, why retrieval depends
on the state of the brain, and why early neurodegeneration can impair access to memories that are still stored. They also show
a connectome at work as an experimental system. Every state variable of every cell can be recorded and moved between brain
states, the model replays bit for bit, and every experiment was declared before it ran. The Discussion turns the findings into
six predictions that can be tested in flies with existing genetic tools.

The core and extended simulation records, trained models, and result tables are deposited at Zenodo
(\href{https://doi.org/10.5281/zenodo.23629827}{10.5281/zenodo.23629827}; \href{https://doi.org/10.5281/zenodo.23629835}{10.5281/zenodo.23629835}), with restricted access during peer review. The lead contact can provide access to editors and reviewers. The code is publicly available at
\url{https://github.com/Hanziwww/fly-connectome-silent-memory}. The authors have approved posting a preprint on bioRxiv under the CC BY-NC-ND 4.0 license.
This manuscript is not under consideration by any other journal. All authors have approved the final manuscript and agree to its submission.

Thank you for considering this work.

Sincerely,

Chunjie Tian, on behalf of all authors\\
Corresponding author and lead contact\\
Department of Otorhinolaryngology, Affiliated Hospital of Guangdong Medical University\\
Zhanjiang 524000, Guangdong, China\\
\href{mailto:tianchunjie@gdmu.edu.cn}{tianchunjie@gdmu.edu.cn}
\end{document}
